\documentclass[a3paper,landscape]{article}
\usepackage[a3paper,landscape,margin=1.5cm]{geometry}
\usepackage{fontspec}
\usepackage{tikz}
\usetikzlibrary{calendar}
\usepackage[dutch]{babel}
\usepackage{etoolbox}

% Gebruik een modern schreefloos lettertype
\setmainfont{Latin Modern Sans}

% --- STAP 1: Lua-script inladen ---
\directlua{dofile("readbdays.lua")}

% --- STAP 2: Weergave macro ---
\newcommand{\showbirthdays}[1]{%
  \edef\currentbdays{\directlua{tex.sprint(get_bdays(\number\pgfcalendarcurrentday, \number\pgfcalendarcurrentmonth))}}%
  \ifdefempty{\currentbdays}{}{%
    \node[anchor=north west, xshift=0.3cm, yshift=-0.5cm, font=\large\bfseries, align=left, text=black] at (0,0) {\currentbdays};
  }%
}

\begin{document}
\pagestyle{empty}

\def\jaar{2026}
\def\breedte{5.5} % Maten in cm zonder eenheid voor Lua
\def\hoogte{4.0}

\foreach \m in {1,...,12}{
    \begin{center}
    \begin{tikzpicture}
        % 1. Teken het grid (7 kolommen, tot 6 rijen)
        \foreach \row in {0,...,5} {
            \foreach \col in {0,...,6} {
                \draw[thick] (\col*\breedte, -\row*\hoogte) rectangle ++(\breedte, -\hoogte);
            }
        }

        % 2. De kalender voor de huidige maand
        \calendar (mycal) [
            dates=\jaar-\m-01 to \jaar-\m-last,
            week list,
            day xshift=\breedte cm,
            day yshift=\hoogte cm,
            day code={
                \node[anchor=north east, xshift=-0.2cm, yshift=-0.2cm, font=\huge\bfseries] at (\breedte, 0) {\tikzdaytext};
                \showbirthdays{\m}
                \ifdate{Monday}{
                    \node[anchor=north west, xshift=0.2cm, yshift=-0.2cm, font=\small\bfseries, text=gray] at (0,0) {W\directlua{tex.sprint(get_week_num(\jaar, \m, \number\pgfcalendarcurrentday))}};
                }{}
            }
        ];

        % 3. De opvulling via Lua (vorige/volgende maand in grijs)
        \directlua{fill_gaps(\jaar, \m, \breedte, \hoogte)}

        % 4. Maandnaam bovenaan
        \node[above, font=\fontsize{50}{60}\selectfont\bfseries, yshift=2.5cm] 
            at (3.5*\breedte, 0) {%
            \ifcase\m\or Januari\or Februari\or Maart\or April\or Mei\or Juni\or
            Juli\or Augustus\or September\or Oktober\or November\or December\fi\ \jaar%
        };

        % 5. Dagnamen
        \foreach \d [count=\i] in {Maandag, Dinsdag, Woensdag, Donderdag, Vrijdag, Zaterdag, Zondag} {
            \node[anchor=south, minimum width=\breedte cm, font=\LARGE\bfseries] 
                at (\i*\breedte - 0.5*\breedte, 0.5) {\d};
        }
    \end{tikzpicture}
    \end{center}
    \newpage
}
\end{document}
